% Report fragment: requires amsmath, amssymb, amsthm, booktabs, longtable.
\section{PA194QNS: the certified fallback construction}

The machine PA194QNS composes the supplied Q axiom emitter, the supplied
no-push modus-ponens interface, and the supplied SUMS implementation of the
original sum-based Shift guard. None of those components is claimed to be
new individually. The new object is their exact 194-control composition,
with complete source and tape certificates and a finite-trace equivalence
to the pinned PA194QN construction. The equivalence is an ordinary
mathematical proof formalizable in PA; no new Lean or Rocq kernel endpoint
is claimed.

\subsection{Literal identity and state budget}

The tape alphabet is $\{0,1\}$, the tape extends in both directions, and
every cell is initially zero. The initial head coordinate is zero. The
working controls are $0,\ldots,193$, with initial control zero and separate
terminal halt label $194$. Every one of the 388 actions writes a bit and
moves exactly one cell left or right. The dense table's SHA-256 is
\begin{quote}\ttfamily\small
7eb36aea69e5de46cc47fa7aab0191f191657ab500dafb808dc3c394640591b8
\end{quote}
The native-label table starts at 74 and halts at 254. Its SHA-256 is
\begin{quote}\ttfamily\small
11613b0b12395fd3132bfe5038c47297685adf337459b138879bb8cdf790daba
\end{quote}
The delivered bijection maps the complete native table to the dense table;
all 194 working controls are graph-reachable.

The unfolded parent has 135 fixed controls and 67 AX readers. The seven
Q-emitter folds and the added comparator fold are simultaneously compatible:
\[
 135+67-8=194.
\]
The 32 local history/read cases are checked against the final simultaneous
quotient. The arithmetic-box rows are unchanged. The separate literal
symbolic checker derives its four list-sweep lemmas from the actual table,
rather than calling the first checker's inherited arithmetic-box oracle.

\subsection{Exact patch relative to QN}

The frozen parent PA194QN has dense digest
\begin{quote}\small\nolinkurl{d1cf0f6e7a3c8f999d045694efc107dd1f9a03d45dd6c15e3d197ff7754c7231}.\end{quote}
The builder reconstructs its entire native and dense tables and source
instruction map before applying the new transformation. The final native
patch is exactly the following; every other row is identical to QN.
\begin{center}
\begin{tabular}{rrll}
\toprule Control&Read&QN action&QNS action\\\midrule
28&0&$(0,R,41)$&$(0,R,330)$\\
38&1&$(1,R,20)$&$(1,R,30)$\\
41&0&$(0,R,24)$&deleted\\
41&1&$(0,R,25)$&deleted\\
70&0&$(0,R,41)$&$(0,R,330)$\\
310&1&$(1,R,62)$&$(1,R,330)$\\
330&0&absent&$(0,L,31)$\\
330&1&absent&$(1,L,46)$\\\bottomrule
\end{tabular}
\end{center}
The no-push action $(164,1)\mapsto(1,R,114)$ is already in QN and is
retained. This avoids claiming it again as a new row change.

Before quotienting, added control 330 has actions $(0,R,63)$ and
$(1,L,73)$. Incoming edges to old control 41 write zero and move right;
incoming edges to raw 330 write one and move right. The immediately
preceding cell therefore distinguishes the two histories. Folding
$41\to330$ gives the displayed final actions and the four three-step
local traces
\[
330,31,21,24;\quad 330,46,61,25;\quad
330,31,47,63;\quad 330,46,62,73.
\]
They cover histories zero/one and current bits zero/one, preserve the
history cell and realize precisely the raw writes, moves and target
controls. The other seven folds retain their full certificates in the
presence of this additional fold.

\subsection{An exact source equivalence}

Both programs use logical registers
$(N,T,W,A,D,H,X,Y,L)$, indexed from zero. Their physical allocation is
$(6,2,0,4,3,7,8,5,1)$. QN has 3,654 primitive source instructions; QNS
has 3,651. Their instructions numbered $0$ through $3648$ are literally
equal, including every resolved target. Only the final five-instruction
tail is replaced by two instructions:
\begin{center}
\begin{tabular}{p{.08\linewidth}p{.43\linewidth}p{.39\linewidth}}
\toprule Site&QN&QNS\\\midrule
3649&decrement $X$; zero $\to3652$, positive $\to3650$&increment $H$; $\to3650$\\
3650&increment $H$; $\to3651$&jump $0$\\
3651&jump $3649$&absent\\
3652&increment $X$; $\to3653$&absent\\
3653&jump $0$&absent\\\bottomrule
\end{tabular}
\end{center}
The only edge from the common prefix into either tail is the zero branch
of instruction 3645, which decrements $H$. Thus $H=0$ on every tail entry
arising from the designated initial source instruction 25. This is checked
by a separate structural certificate that imports neither compiler nor
physical executor.

The common instructions $0,\ldots,10$ form a destructive equality test:
on input $(H,X)=(u,v)$ they halt exactly when $u=v$; otherwise they
terminate at instruction 11 with $H=X=0$ and all other registers
unchanged. Lockstep decrements detect simultaneous exhaustion; the
remaining loops clear an unmatched counter. Induction on $u+v$ proves
termination and the asserted result in PA.

At entry 3649 write $X=x$; the incoming invariant gives $H=0$. QN's
finite tail transfers $X$ to $H$ and then increments $X$, yielding
comparator operands $(x,1)$. QNS simply increments $H$, yielding
$(1,x)$. They therefore both halt exactly when $x=1$, and otherwise
return to the identical complete store at instruction 11. The transfer
loop terminates by induction on its initial value $x$.

More explicitly, the common Shift prefix at 3642 transfers incoming
$T=t$ into $X$, then incoming $H=h$ into both $T$ and $X$. Starting with
arbitrary $(T,H,X)=(t,h,x)$, it reaches the tail with
\[
 (T,H,X)=(h,0,x+t+h).
\]
Both complete Shift blocks consequently halt exactly when $x+t+h=1$;
otherwise they leave $(T,H,X)=(h,0,0)$ and preserve every other register.
At command boundaries $X=0$, recovering precisely the original guard
$t+h=1$. No cached-head-only guard is being substituted.

For finite-trace transport, cut a source execution whenever it enters
3649 and follow that local block through the comparator to instruction 11
or halt. Replace each such terminating block by its counterpart. Outside
these blocks the resolved instruction graphs are identical. The sole-entry
check supplies $H=0$ at every replacement, and every block has positive
finite duration. Induction on the number of replacements gives
primitive-recursive translations of finite halting traces in both
directions, from any complete register store at source 25. PA proves
the correctness of these translations. In particular,
\[
 \mathrm{PA}\vdash
 \operatorname{Halt}(\mathrm{source}_{\rm QNS},25)
 \leftrightarrow
 \operatorname{Halt}(\mathrm{source}_{\rm QN},25).
\]

\subsection{From the source to consistency}

The complete blank computation of QNS reaches the clean source-25 frame
after 58,004 physical steps, at native control 84 and head coordinate
$-2$. The logical registers are $(2,0,0,0,0,0,0,0,0)$, followed by 26
one-valued spare registers. Both simulators compare the entire finite
tape support and zero exterior, not a tape window. QN reaches that
source frame after 58,154 steps.

The universal source-frame certificates cover 4,751 primitive outcomes,
including both branches of every decrement, arbitrary natural primary
registers, an arbitrary affine header gap, and the stated complete spare
list interface. Together with positive progress and deterministic tape
execution, they identify blank halting with halting of the corresponding
source from its initial frame. The new source equivalence therefore
transports the supplied QN paper calibration:
\[
 \mathrm{PA}\vdash
 \operatorname{Halt}(\mathrm{PA194QNS},\mathrm{blank})
 \leftrightarrow \exists e\,
       \operatorname{Prf}_{\mathrm{PA}}(e,0=1),
\]
and hence
\[
 \mathrm{PA}\vdash
 \neg\operatorname{Halt}(\mathrm{PA194QNS},\mathrm{blank})
 \leftrightarrow\operatorname{Con}(\mathrm{PA}).
\]
The inherited QN calibration retains the separate diagonal-head calculus,
the K/S exchange derivation for Q's curried-premise change, and the
reserved-zero proof compiler for no-push MP. They are attributed in the
reference fragments and are not proved again merely by executing the
physical certificates. This transport needs no assumption that PA is
consistent. Such an assumption is used only to conclude actual
nonhalting from the biconditional.

\subsection{Material difference and its limits}

QNS changes the executing Shift protocol of QN, eliminating a counter
transfer on every visit while retaining its arbitrary-store behavior.
It changes the Q axiom emitter relative to the supplied sum-symmetric
NS machine. These are recurring logical and runtime operations, not
unreachable-row edits. Its complete action graph is absent from all 39
current 194-state graph classes at the pinned revision, even after
state renaming, spanning-root change, reflection and bit interchange.
The component techniques are inherited, and no sub-194 construction or
global minimality result is claimed for this fallback.
